\chapter{Glossary}
\label{appendix:glossary}

\section{Technical Terms}

\begin{description}
    \item[API (Application Programming Interface)] A set of protocols and tools for building software applications.
    
    \item[ARPU (Average Revenue Per User)] The average revenue generated per user over a specific period.
    
    \item[CAC (Customer Acquisition Cost)] The cost of acquiring a new customer.
    
    \item[COGS (Cost of Goods Sold)] Direct costs attributable to the production of goods sold.
    
    \item[EBITDA] Earnings Before Interest, Taxes, Depreciation, and Amortization.
    
    \item[JWT (JSON Web Token)] A compact, URL-safe means of representing claims to be transferred between two parties.
    
    \item[LTV (Lifetime Value)] The predicted net profit attributed to the entire future relationship with a customer.
    
    \item[Microservices] An architectural style that structures an application as a collection of loosely coupled services.
    
    \item[MRR (Monthly Recurring Revenue)] The predictable revenue that a company expects to receive every month.
    
    \item[OAuth] An open standard for access delegation, commonly used for token-based authentication.
    
    \item[RBAC (Role-Based Access Control)] A method of regulating access to resources based on the roles of individual users.
    
    \item[REST (Representational State Transfer)] An architectural style for designing networked applications.
    
    \item[SaaS (Software as a Service)] A software licensing and delivery model in which software is licensed on a subscription basis.
    
    \item[SOC 2] A compliance framework for service organizations to manage customer data.
    
    \item[TAM (Total Addressable Market)] The total market demand for a product or service.
\end{description}
